% =====================================================================
% bucket_table_ci.tex
%
% Bucket experiment: normalized probability mass on REAL vs HALL object
% candidates, WITH 95% bootstrap confidence intervals (B = 10,000, resampling
% the 500 images). Point estimates are unchanged from the original table; only
% the intervals are added.
%
% Requires in the preamble:  \usepackage{booktabs}  \usepackage{multirow}
%                            \usepackage{float}      % for the [H] placement
% Table number is automatic; refer to it with \ref{tab:bucket_v3_combined}.
%
% Generated by bucket_ci.py (image-level cluster bootstrap).
% =====================================================================

\begin{table}[H]
\centering
\caption{Normalized probability mass assigned to present (REAL) and
hallucinated (HALL) object candidates, with 95\% bootstrap confidence intervals
($B=10{,}000$, resampling the 500 images). HALL $=1-$REAL by construction, so its
interval is the mirror of REAL's and is not an independent estimate.}
\label{tab:bucket_v3_combined}
\setlength{\tabcolsep}{4pt}
\renewcommand{\arraystretch}{1.15}
\begin{tabular}{lllcc}
\toprule
\textbf{Model} & \textbf{Method} & \textbf{Branch}
& \textbf{REAL} & \textbf{HALL} \\
\midrule
\multirow{6}{*}{LLaVA-v1.5-7B}
& \multirow{3}{*}{VCD} & Expert  & $0.84027$ \tiny{[0.82599, 0.85478]} & $0.15973$ \tiny{[0.14522, 0.17401]} \\
& & Amateur & $0.84007$ \tiny{[0.82540, 0.85523]} & $0.15993$ \tiny{[0.14477, 0.17460]} \\
& & CD      & $0.83964$ \tiny{[0.82509, 0.85426]} & $0.16036$ \tiny{[0.14574, 0.17491]} \\
\cmidrule(lr){2-5}
& \multirow{3}{*}{SID} & Expert  & $0.84459$ \tiny{[0.82988, 0.85897]} & $0.15541$ \tiny{[0.14103, 0.17012]} \\
& & Amateur & $0.84021$ \tiny{[0.82519, 0.85530]} & $0.15979$ \tiny{[0.14470, 0.17481]} \\
& & CD      & $0.84381$ \tiny{[0.82911, 0.85829]} & $0.15619$ \tiny{[0.14171, 0.17089]} \\
\midrule
\multirow{6}{*}{Qwen2.5-VL-7B-Instruct}
& \multirow{3}{*}{VCD} & Expert  & $0.89641$ \tiny{[0.88077, 0.91116]} & $0.10359$ \tiny{[0.08884, 0.11923]} \\
& & Amateur & $0.90854$ \tiny{[0.89372, 0.92246]} & $0.09146$ \tiny{[0.07754, 0.10628]} \\
& & CD      & $0.88893$ \tiny{[0.87239, 0.90440]} & $0.11107$ \tiny{[0.09560, 0.12761]} \\
\cmidrule(lr){2-5}
& \multirow{3}{*}{SID} & Expert  & $0.89615$ \tiny{[0.88190, 0.90989]} & $0.10385$ \tiny{[0.09011, 0.11810]} \\
& & Amateur & $0.91472$ \tiny{[0.90091, 0.92760]} & $0.08528$ \tiny{[0.07240, 0.09909]} \\
& & CD      & $0.89083$ \tiny{[0.87613, 0.90531]} & $0.10917$ \tiny{[0.09469, 0.12387]} \\
\bottomrule
\end{tabular}
\end{table}


% =====================================================================
% Method + interpretation to place after the table.
% =====================================================================

\paragraph{How the confidence intervals were computed.} Each reported value is
the mean, over all object-emitting steps, of a per-step ratio. At a given step we
take the branch's top-5 candidate tokens together with their softmax
probabilities (softmax taken over that branch's stored top-30 logits), keep only
the tokens that are object words, and form
$\mathrm{REAL} = \sum P_{\mathrm{real}} / \left[\sum P_{\mathrm{real}} + \sum P_{\mathrm{hall}}\right]$,
with $\mathrm{HALL} = 1 - \mathrm{REAL}$. Intervals come from a cluster bootstrap
whose resampling unit is the \emph{image}: for each of the 500 images we store the
sum of the per-step ratios and the number of contributing steps, draw $B=10{,}000$
resamples of the 500 images with replacement, and recompute the statistic on each
resample as a ratio of sums (total ratio mass divided by total steps). The
reported interval is the 2.5th to 97.5th percentile of the $10{,}000$ replicates.
We resample images rather than steps because steps inside one caption share the
same picture and the same ground-truth object set and are therefore correlated;
resampling steps independently would understate the uncertainty. For the paired
differences below, a single image resample is applied to all three branches
within each replicate, so the shared image-sampling variance cancels.

\paragraph{What the intervals show.} The intervals are narrow (about $\pm 0.015$
on each estimate), and within every model and method the Expert, Amateur, and CD
intervals overlap almost completely: the split of object probability mass between
present and hallucinated candidates is essentially unchanged by the contrastive
step. Because the three branches are evaluated on the same steps, we test the
branch differences with the paired bootstrap rather than by eyeballing overlapping
marginal intervals. On LLaVA the change contrastive decoding makes is
indistinguishable from zero (CD $-$ Expert $= -0.0006$, 95\% CI
$[-0.0033, +0.0022]$ for VCD; $-0.0008$, $[-0.0034, +0.0019]$ for SID), and the
interval is tight enough to bound any true shift below roughly $0.003$. That is a
bounded null, i.e.\ evidence that contrastive decoding does not meaningfully
redistribute object mass, not merely an absence of evidence. On Qwen the change is
statistically significant and in the unintended direction: contrastive decoding
\emph{lowers} the mass placed on present objects by $0.0075$ (VCD, 95\% CI
$[-0.0107, -0.0044]$, $p \approx 0.0001$) and by $0.0053$ (SID,
$[-0.0084, -0.0024]$, $p \approx 0.0002$). The mechanism is visible in the
Amateur row: on Qwen the degraded branch is significantly \emph{more}
concentrated on present objects than the expert ($+0.0121$ under VCD, $+0.0186$
under SID), so subtracting it necessarily removes mass from correct objects.

Both outcomes support our claim. Where contrastive decoding looks harmless
(LLaVA) it provably changes nothing about which objects receive probability mass;
where it does change something (Qwen) it moves mass \emph{towards} hallucinated
objects. Note that the LLaVA result should be stated as ``no measurable change''
rather than ``a shift in the wrong direction'': the small negative point
differences there are not statistically distinguishable from zero.

\paragraph{Why the paired test rather than the overlapping intervals.} The
intervals in Table~\ref{tab:bucket_v3_combined} are \emph{marginal}: each
describes the uncertainty of one branch taken on its own. Almost all of that
width is image-sampling variance, and because Expert, Amateur and CD are computed
on exactly the same object steps of exactly the same images, that component is
\emph{shared} between them and cancels when the branches are compared against
each other. Deciding whether two branches differ by checking whether their
marginal intervals overlap therefore discards the pairing and is badly
underpowered. Table~\ref{tab:bucket_paired_diff} instead draws one image resample
per replicate and recomputes all three branches on that same resample, so every
replicate yields a difference from which the shared variance has already
cancelled; the reported interval is the uncertainty of the difference itself.

This distinction changes the conclusion, so it is not a formality. On Qwen under
VCD the marginal intervals for Expert ($[0.88077, 0.91116]$) and CD
($[0.87239, 0.90440]$) overlap across most of their range, which would invite the
reading that the two branches are indistinguishable. The paired difference is
$-0.00747$ with a 95\% interval of $[-0.01074, -0.00439]$, excluding zero at
$p \approx 0.0001$: comparing the marginal intervals alone would have missed a
real effect. Conversely, on LLaVA the paired interval contains zero and is narrow,
and it is that narrowness which licenses the stronger claim that contrastive
decoding produces no measurable shift there, rather than merely an unresolved one.

% (optional) companion table: the paired branch differences, the correct test
\begin{table}[H]
\centering
\caption{Paired bootstrap differences in REAL mass ($B=10{,}000$, the same image
resample applied to all branches within each replicate). An interval excluding
zero indicates a genuine branch difference.}
\label{tab:bucket_paired_diff}
\setlength{\tabcolsep}{6pt}
\renewcommand{\arraystretch}{1.15}
\begin{tabular}{llcc}
\toprule
\textbf{Model} & \textbf{Method} & \textbf{CD $-$ Expert} & \textbf{Amateur $-$ Expert} \\
\midrule
LLaVA-v1.5-7B & VCD & $-0.00064$ \tiny{[$-0.00331$, $+0.00215$]} & $-0.00021$ \tiny{[$-0.00439$, $+0.00394$]} \\
              & SID & $-0.00078$ \tiny{[$-0.00344$, $+0.00188$]} & $-0.00438$ \tiny{[$-0.00813$, $-0.00060$]} \\
\midrule
Qwen2.5-VL-7B-Instruct & VCD & $-0.00747$ \tiny{[$-0.01074$, $-0.00439$]} & $+0.01213$ \tiny{[$+0.00682$, $+0.01758$]} \\
                       & SID & $-0.00532$ \tiny{[$-0.00841$, $-0.00235$]} & $+0.01857$ \tiny{[$+0.01194$, $+0.02535$]} \\
\bottomrule
\end{tabular}
\end{table}
